% 编译方式: xelatex SL_stability_moment_jump.tex
\documentclass[UTF8]{ctexart}
\usepackage{amsmath,amssymb,amsthm}
\usepackage[margin=2.5cm]{geometry}
\usepackage[hyphens]{url}
\usepackage[hidelinks]{hyperref}
\usepackage{booktabs}
\emergencystretch 3em
\newtheorem{theorem}{定理}[section]
\newtheorem{lemma}[theorem]{引理}
\newtheorem{corollary}[theorem]{推论}
\newtheorem{remark}[theorem]{注}
\newtheorem{definition}[theorem]{定义}
\newtheorem{proposition}[theorem]{命题}

\title{矩跳跃的条件稳定性\\ \large 一般递推下界、乘积模型与基扰动反例}
\author{研究证明文档 (项目: BVE research)}
\date{2026-09-20, 第三轮审计修订}

\begin{document}
\maketitle
\begin{abstract}
一般二阶矩递推给出的是乘积下界; 系数差的对数总和控制该下界,
不能单独决定实际解的上界或完备性. 本文保留定量增长引理和分别作用于
奇、偶递推的完备性充分条件, 给出以实际解为对象的对角级数充要判据.
精确乘积及显式临界指数只在明确的 $B_m=B'_m=0$ 模型中使用.
对基系数扰动同时计算 $A_m$ 与 $B_m$: 有界甚至趋零的扰动可以破坏
$L^2$ 中形式像族的完备性, 也通常离开真正的 Krein 算子域.
文末给出非零正交元反例与一个具有全指标假设的扰动充分条件.
本修订撤回旧版的一般增长分类和无条件有界基扰动稳健性,
不撤回原始 Krein 族已另行证明的低阶结论.
\end{abstract}
\tableofcontents

\section{对象、空间与矩递推}

全文的抽象论证在实 Hilbert 空间 $H$ 中进行. 假设
\begin{description}
\item[(H1)] 多项式空间 $\Pi\subset H$ 且在 $H$ 中稠密.
\item[(H2)] 矩使用同一内积 $M_k=(w,x^k)_H$; 因 $x^k\in H$,
它们是关于 $w$ 的连续线性泛函.
\item[(H3)] 存在 $C_H>0,\ \beta\ge0$, 对每个 $k\ge0$,
$\|x^k\|_H\le C_H(k+1)^\beta$.
\end{description}
使用 $k+1$ 也包含常数模态, 不在 $k=0$ 人为产生零范数界.
给定 $c_0>0$, 考虑\emph{缺少指标 2,3} 的族
\begin{equation}\label{eq:family}
\begin{aligned}
q_0&=c_0,\qquad q_1=c_0x,\\
q_{2m}&=c_0x^{2m}-A_mx^{2m-2}+B_mx^{2m-4},\\
q_{2m+1}&=c_0x^{2m+1}-A'_mx^{2m-1}+B'_mx^{2m-3},
\qquad m\ge2 .
\end{aligned}
\end{equation}
基本解由
\begin{equation}\label{eq:recurrence}
u_0=0,\quad u_1=1,\quad
c_0u_m=A_mu_{m-1}-B_mu_{m-2}\quad(m\ge2)
\end{equation}
定义, 奇侧 $u'_m$ 使用 $A'_m,B'_m$.
若 $w$ 正交于 \eqref{eq:family}, 则
\begin{equation}\label{eq:moments}
M_0=M_1=0,\qquad M_{2m}=M_2u_m,\quad M_{2m+1}=M_3u'_m
\quad(m\ge0).
\end{equation}
最后两个等式由初值及递推归纳得到, 无需假定 $M_2$ 或 $M_3$ 非零.

这里的 $L_{\rm poly}=c-D^2$ 是多项式上的形式微分表达式.
只有先核对 $p\in D(K_c)$, 才能将 $L_{\rm poly}p$ 称为真正的 $K_cp$.
抽象前提 (H1)--(H3) 也不能直接套到不含全部单项式的算子幂域.
原稀疏族的 $H^2/H^3$ 完备性及其 $0\le s\le3$ 谱截断推论另见相关文稿;
原同族在 $s\ge7/2$ 的断言已撤回, $3<s<7/2$ 的稠密性仍未判定.

\section{一般系数的增长下界与充分性}

\begin{theorem}[定量增长下界]\label{thm:growth}
设对每个 $m\ge2$ 有 $B_m\ge0$ 且 $A_m-B_m\ge c_0>0$.
令 $\varepsilon_m=(A_m-B_m-c_0)/c_0\ge0$.
则 \eqref{eq:recurrence} 的解满足
\begin{equation}\label{eq:decomposition}
u_m=(1+\varepsilon_m)u_{m-1}
+\frac{B_m}{c_0}(u_{m-1}-u_{m-2}),
\end{equation}
并且单调不减, $u_m\ge1$ ($m\ge1$), 以及
\begin{equation}\label{eq:product-lower}
u_m\ge \prod_{k=2}^m\frac{A_k-B_k}{c_0}
=\prod_{k=2}^m(1+\varepsilon_k).
\end{equation}
\end{theorem}
\begin{proof}
等式 \eqref{eq:decomposition} 是原递推的恒等改写.
由 $u_1=1\ge u_0=0$, 若 $u_{m-1}\ge u_{m-2}\ge0$,
则右端第二项非负, 第一项不小于 $u_{m-1}$.
归纳得单调性, 再迭代 $u_m\ge(1+\varepsilon_m)u_{m-1}$
得到 \eqref{eq:product-lower}. 空乘积对应 $u_1=1$.
\end{proof}

\begin{remark}[下界与等式]
$B_m=0$ 时相应一步为乘积等式; 一般 $B_m>0$ 时不得丢弃
\eqref{eq:decomposition} 的第二项后仍声称等式.
$u_m\ge (A_m/c_0)u_{m-1}$ 也一般不成立:
原 Krein 系数在 $c=3$ 时 $u_2=6,u_3=63,u_4=1180<1260=(A_4/c)u_3$.
以上全指标条件是足够条件, 不宣称刻画所有单调递推.
仅有 $A_m\ge B_m$ 不够, 例如 $c_0=1,A_m=B_m=1$
给出 $0,1,1,0,-1,\ldots$.
\end{remark}

\begin{theorem}[奇偶两侧的稳定性充分条件]\label{thm:stability}
在 (H1)--(H3) 与 \eqref{eq:family} 的设置下, 假设两侧对所有 $m\ge2$
都满足定理 \ref{thm:growth} 的系数条件, 并且
\[
\frac{\sum_{k=2}^m\min(\varepsilon_k,1)}{\log m}\longrightarrow\infty,
\qquad
\frac{\sum_{k=2}^m\min(\varepsilon'_k,1)}{\log m}\longrightarrow\infty .
\]
则族 \eqref{eq:family} 在 $H$ 中完备.
\end{theorem}
\begin{proof}
对 $t\ge0$, $\log(1+t)\ge(\log2)\min(t,1)$.
定理 \ref{thm:growth} 于是给出两侧 $\log u_m/\log m\to\infty$
和 $\log u'_m/\log m\to\infty$.
若 $w$ 正交于该族, \eqref{eq:moments} 与 Cauchy--Schwarz 给出
\[
|M_2|u_m=|M_{2m}|\le \|w\|_H C_H(2m+1)^\beta .
\]
超多项式增长迫使 $M_2=0$, 奇侧同理迫使 $M_3=0$.
所有矩均为零, (H1) 给出 $w=0$, 故正交补为零.
\end{proof}

\begin{theorem}[$S$-条件: 更弱的充分条件]\label{thm:Sthr}
保留上一定理的空间、跳变结构和两侧全指标系数条件, 令
\[
S(m)=\sum_{k=2}^m\log(1+\varepsilon_k),\qquad
S'(m)=\sum_{k=2}^m\log(1+\varepsilon'_k).
\]
若 $S(m)/\log m\to\infty$ 且 $S'(m)/\log m\to\infty$,
则该族在 $H$ 中完备.
\end{theorem}
\begin{proof}
只需使用 $\log u_m\ge S(m)$、$\log u'_m\ge S'(m)$,
再应用上一定理中的矩界论证. 此处不使用等式 $\log u_m=S(m)$.
\end{proof}

\begin{proposition}[稀疏大跳的充分性]\label{prop:sparse}
在定理 \ref{thm:Sthr} 的设置下, 取两侧
$\varepsilon_k=\varepsilon'_k=e^k$ 当 $k=k_j=2^{2^j}$ ($j\ge1$),
其余为零. 则两侧满足 $S$-条件, 而
$\sum_{k\le m}\min(\varepsilon_k,1)=O(\log\log m)=o(\log m)$.
\end{proposition}
\begin{proof}
若 $k_J\le m<k_{J+1}=k_J^2$, 则最大的跳点 $k_J>\sqrt m$.
所以 $S(m)\ge\log(1+e^{k_J})\ge k_J>\sqrt m$.
跳点个数为 $O(\log\log m)$, 两侧相同.
这证明 $S$-条件确实比截断增量条件覆盖更多情形.
\end{proof}

\section{对角空间: 实际解判据与纯乘积模型}

对任意有限实数 $\beta$, 定义 $H_\beta$ 为多项式在内积
$(x^j,x^k)=\delta_{jk}(k+1)^{2\beta}$ 下的完备化.
这是系数对角 Hilbert 空间:
\[
w=\sum_{k\ge0}a_kx^k,\qquad
\|w\|_\beta^2=\sum_{k\ge0}|a_k|^2(k+1)^{2\beta}.
\]
它不同于 Krein 左定空间 $H^s$, 不以相同字母的外形混用两种定义.

\begin{theorem}[以实际递推解为对象的精确判据]\label{thm:exact}
对 \eqref{eq:family} 及其实际基本解 $u,u'$, 在 $H_\beta$ 中完备当且仅当
\begin{equation}\label{eq:series}
\sum_{m\ge1}\frac{|u_m|^2}{(2m+1)^{2\beta}}=\infty
\quad\hbox{且}\quad
\sum_{m\ge1}\frac{|u'_m|^2}{(2m+2)^{2\beta}}=\infty .
\end{equation}
本判据只需 $c_0>0$ 和所列跳变结构, 不需要递推解单调.
\end{theorem}
\begin{proof}
对角性给出 $M_k=a_k(k+1)^{2\beta}$.
正交条件等价于 \eqref{eq:moments}.
若两个级数都发散, 有限范数要求 $M_2=M_3=0$, 故 $w=0$.
若偶侧级数收敛, 则
\[
w=\sum_{m\ge1}u_m(2m+1)^{-2\beta}x^{2m}
\]
属于 $H_\beta$, 因 $u_1=1$ 而非零, 且逐一满足全部正交条件;
奇侧收敛时用 $u'_m(2m+2)^{-2\beta}x^{2m+1}$ 同理.
\end{proof}

\begin{corollary}[实际解上界给出的不完备性]\label{cor:di}
若某一个奇偶基本解满足 $|u_m|\le K(m+1)^D$ (或对 $u'_m$ 有此界),
则在 $\beta>D+1/2$ 时族 \eqref{eq:family} 在 $H_\beta$ 中不完备.
\end{corollary}
\begin{proof}
相应级数受常数乘 $\sum m^{2D-2\beta}$ 控制而收敛.
这个上界是实际解的额外假设, 不能由 $S(m)=O(\log m)$ 单独推出.
\end{proof}

\begin{theorem}[$B=0$ 的 $C/k$ 乘积模型]\label{thm:sharp}
令 $C>0$, 并对所有 $k\ge2$ 明确取
$B_k=B'_k=0$, $A_k=A'_k=c_0(1+C/k)$.
则
\begin{equation}\label{eq:gamma}
u_m=u'_m=\prod_{k=2}^m(1+C/k)
=\frac{\Gamma(m+1+C)}{\Gamma(2+C)\Gamma(m+1)}
=\frac{m^C}{\Gamma(2+C)}(1+o(1)).
\end{equation}
因此 $\beta>C+1/2$ 时有定理 \ref{thm:exact} 中的非零正交元.
\end{theorem}
\begin{proof}
$B=B'=0$ 将递推化为一阶乘积. Gamma 递推公式给出有限恒等式,
Stirling 公式给出末尾渐近. 对应范数级数与 $\sum m^{2C-2\beta}$
同敛散, 因而所给正交元在指定范围内范数有限.
起始指标为 $2$, 归一化常数是 $\Gamma(2+C)$; 例如 $C=2$ 时
$u_m=(m+1)(m+2)/6$, 而不是以 $1/2$ 为首项常数.
\end{proof}

\begin{theorem}[$C/k$ 模型的临界指数]\label{thm:model}
在定理 \ref{thm:sharp} 的全部假设下, 在 $H_\beta$ 中完备
当且仅当 $\beta\le C+1/2$.
\end{theorem}
\begin{proof}
将 \eqref{eq:gamma} 代入奇偶两个级数.
在等号 $\beta=C+1/2$ 处均与调和级数同敛散, 因而发散.
\end{proof}

\begin{theorem}[$B=0$ 的对数乘积模型]\label{thm:line}
对每个 $k\ge2$ 取 $B_k=B'_k=0$ 且
$A_k=A'_k=c_0(1+C/(k\log k))$, $C>0$.
则
\[
\log u_m=\log u'_m=C\log\log m+O(1),\qquad
u_m=u'_m\asymp(\log m)^C,
\]
且在 $H_\beta$ 中完备当且仅当 $\beta\le1/2$.
\end{theorem}
\begin{proof}
写 $t_k=C/(k\log k)$. 由
$0\le t_k-\log(1+t_k)\le t_k^2/2$ 及 $\sum t_k^2<\infty$,
乘积对数与 $\sum t_k$ 相差有界量.
积分比较又给出 $\sum_{k=2}^m1/(k\log k)=\log\log m+O(1)$.
所以这里得到的是有界误差, 足以推出所写的双边 $\asymp$;
仅有对数的渐近比值并不足以推出它.
判据 \eqref{eq:series} 化为
$\sum m^{-2\beta}(\log m)^{2C}$ 的敛散性.
$\beta=1/2$ 时它大于常数乘调和级数;
$\beta>1/2$ 时用对数幂小于任意正幂并作比较; $\beta<1/2$ 时发散.
\end{proof}

\begin{proposition}[相同系数差不决定实际增长]\label{prop:counter-recurrence}
取 $c_0=1$. 两侧均取 $A_m=3,B_m=2$ 时,
$\varepsilon_m=0,S(m)=0$, 却有
$u_m=u'_m=2^m-1$, 故在每个有限 $\beta$ 的 $H_\beta$ 中完备.
两侧均取 $B_m=2,A_m=3+2/m$ 时也有 $u_m,u'_m\ge2^m-1$,
而乘积模型只给出 $(m+1)(m+2)/6$.
\end{proposition}
\begin{proof}
$2^m-1$ 满足初值与常系数递推, 由唯一性得到第一式.
第二种情形的 $d_m=u_m-u_{m-1}$ 满足
$d_m=2d_{m-1}+(2/m)u_{m-1}\ge2d_{m-1}$.
定理 \ref{thm:growth} 保证 $u_{m-1}\ge0$, $d_1=1$,
故 $d_m\ge2^{m-1}$ 及 $u_m\ge2^m-1$.
$m=2$ 时实际值为 $4$, 乘积值为 $2$, 等式在首步已失败.
两侧指数增长使 \eqref{eq:series} 对任意有限 $\beta$ 都发散.
将 $2/m$ 换成任意非负增量, 包括 $C/(m\log m)$,
同一个差分下界仍成立.
\end{proof}

纯乘积模型给出了充分条件不能任意放宽的具体例子,
不构成一般 $B\ge0$ 的完整分类.
此外 $S=O(\log m)$ 与 $S/\log m\to\infty$ 不穷尽所有非负增量序列.
一般非对角空间中的矩可表示性还需要内积/Gram 结构;
仅有单项式范数的下界不能替代对角判据.

\section{Krein 系数与基扰动}

\begin{theorem}[固定正移位与有条件的基扰动]\label{thm:margin}
令 $c>0$. 原偶、奇形式系数分别为
\[
A_m=2m(2m-1)+\frac{cm}{m-1},\quad B_m=2m(2m-3),
\]
\[
A'_m=2m(2m+1)+\frac{cm}{m-1},\quad B'_m=2m(2m-1).
\]
两侧均有 $A_m-B_m=4m+cm/(m-1)\ge4m+c$.
\begin{enumerate}
\item 对任意新的 $c'>0$, 使用与 $c'$ 一致的首项系数和两侧
$A,B$ 后仍满足增长条件; 基本解至少为
$(4/c')^{m-1}m!$, 因而在满足 (H1)--(H3) 的抽象空间中形式像族完备.
\item 对偶侧
$\widetilde p_{2m}=x^{2m}-(m/(m-1)+\delta_m)x^{2m-2}$, 其形式像的真实系数为
\begin{align}
A_m^\delta&=2m(2m-1)+c\left(\frac m{m-1}+\delta_m\right),\label{eq:perturb-a}\\
B_m^\delta&=\left(\frac m{m-1}+\delta_m\right)(2m-2)(2m-3),\label{eq:perturb-b}\\
A_m^\delta-B_m^\delta
&=4m+\frac{cm}{m-1}+\delta_m\{c-(2m-2)(2m-3)\}.\label{eq:perturb-gap}
\end{align}
奇侧以 $\delta'_m$ 扰动相同的基系数, 对应
$B_m^{\prime,\delta}=(m/(m-1)+\delta'_m)(2m-1)(2m-2)$,
差式中的二阶导数因子随之替换.
若\emph{每个} $m\ge2$ 的实际两侧系数都满足
$B^\delta\ge0,A^\delta-B^\delta\ge c$, 且两侧 $S$-条件成立,
则抽象形式像族完备.
一个显式足够条件是对所有 $m\ge2$,
$|\delta_m|,|\delta'_m|\le1/(8m)$.
\end{enumerate}
这些基扰动结论没有宣称扰动多项式属于真正的 $D(K_c)$.
\end{theorem}
\begin{proof}
未扰动的差式直接计算, 将 $c$ 换成 $c'>0$ 后连乘
$(A_k-B_k)/c'\ge4k/c'$ 得第一项.
对扰动必须对两个单项式都求二阶导数, 得
\eqref{eq:perturb-a}--\eqref{eq:perturb-gap}; 条件结论来自定理 \ref{thm:Sthr}.
对显式足够条件, 记 $\alpha_m=m/(m-1)+\delta_m$,
则 $\alpha_m\ge m/(m-1)-1/(8m)>1$.
偶、奇侧的导数因子 $D_m$ 都在 $[0,4m^2]$ 中, 所以
\[
B_m^\delta=\alpha_mD_m\ge0,\qquad
A_m^\delta-B_m^\delta=4m-\delta_mD_m+c\alpha_m
\ge \tfrac72m+c .
\]
两侧对数下界于是超多项式, 给出所述充分条件.
\end{proof}

\begin{remark}[域与有限前缀不可略去]
$c=3,\delta_4=1$ 时 $A_4^\delta=63,B_4^\delta=70$,
真实差为 $-7$, 漏掉 $B$ 的扰动会错误地算成 $23$.
而 $\widetilde p_{2m}$ 为偶函数,
$\widetilde p'_{2m}(1)=-(2m-2)\delta_m$;
真正 Krein 边界条件要求偶函数的此导数为零.
奇侧也有
$\widetilde p'_{2m+1}(1)-\widetilde p_{2m+1}(1)=-(2m-2)\delta'_m$.
非零的这类扰动一般已离开 $D(K_c)$.

仅从某个下标以后满足增长条件不够: 前面的有限递推可能消去增长分量.
下面反例的扰动甚至最终满足上述 $1/(8m)$ 界,
却因全指标条件不成立而不完备.
\end{remark}

\begin{proposition}[有界且趋零的基扰动可以不完备]\label{prop:perturb-counter}
在 $H=L^2(-1,1)$、$c=3$ 的形式多项式设置下,
存在 $-36/7\le\delta_m\le1$ 且
$\delta_m\sim3/(2m^2)$ 的偶侧基扰动, 保留原奇侧族,
使 $q_0=3,q_1=3x$ 及其余形式像组成的族不完备.
\end{proposition}
\begin{proof}
取 $g(x)=P_2(x)=(3x^2-1)/2$, 则
$(g,1)=(g,x)=0$, $\|g\|_2^2=2/5>0$.
直接积分对 $m\ge1$ 给出
\[
\mu_{2m}=(g,x^{2m})=\frac{4m}{(2m+1)(2m+3)},\qquad
\nu_{2m}=(g,L_{\rm poly}x^{2m})
=-\frac{4m(4m^2+2m-9)}{(2m+1)(2m+3)}.
\]
$m=1$ 时 $\nu_2=4/5$, $m\ge2$ 时 $\nu_{2m}<0$, 故均非零.
对 $m\ge2$ 取 $\alpha_m=\nu_{2m}/\nu_{2m-2}$
及 $\delta_m=\alpha_m-m/(m-1)$, 化简为
\[
\delta_m=\frac{6m(2m+5)}
{(m-1)(2m+3)(4m^2-6m-7)} .
\]
有 $\delta_2=-36/7,\delta_3=1$; 对 $m\ge3$, 分母为正, $\delta_m>0$, 且
\[
1-\delta_m=
\frac{(m-3)(2m-1)(4m^2+10m+7)}
{(m-1)(2m+3)(4m^2-6m-7)}\ge0.
\]
最高次项给出 $\delta_m\sim3/(2m^2)$.
按定义 $(g,L_{\rm poly}\widetilde p_{2m})
=\nu_{2m}-\alpha_m\nu_{2m-2}=0$ 对\emph{所有} $m\ge2$ 成立.
原奇多项式及其形式像都是奇函数, 也与偶函数 $g$ 正交;
$q_0,q_1$ 已检验. 因而 $g$ 是整个族的非零正交元.
这里 $\alpha_2=-22/7$, 实际 $B_2^\delta<0$,
不满足增长定理的全指标假设, 没有反驳正确的条件定理.
\end{proof}

\section{验证边界与保留的问题}

本次原始外部报告、旧源文、作者精确检查和独立复核保存在
\path{research/artifacts/proof-audit-round3-20260920/};
实际结果及形式化范围见
\path{reports/proof-audit-round3-20260920/REPORT.md}.
旧稿中的数值列表不能证明一般增长分类或有界基扰动稳定性, 已从当前证明依据中撤除.

活动脚本 \path{scripts/d3_stability_verify.py}、
\path{scripts/d3_stability_verify2.py} 核查实际递推与系数;
\path{scripts/op12_dichotomy_verify.py}、
\path{scripts/op12_threshold_verify.py}、
\path{scripts/op12_sparse_check.py} 只对明确的乘积模型作有限数值诊断.
精确反例的全指标结论由上述代数恒等式和归纳证明承担,
不由有限扫描或部分和的大小承担. 局部 Lean 目标不等于整个 Hilbert 空间完备性定理已形式化.

保留的问题包括: 一般二阶系数的完整增长分类; 非对角空间的矩可表示性;
真实算子域中允许哪些结构保持扰动; 以及变系数微分表达式的闭合机制.
原始 Krein 族的低阶证明、首对比值、MW 和 K1 等支线不因本轮反例而一并撤回.

\begin{thebibliography}{4}
\bibitem{h2proof} 项目文档, \emph{$H^2$ 多项式基解析完备性: 证明文档}.
\bibitem{h3proof} 项目文档, \emph{$H^3$ 多项式完备性: 证明文档}.
\bibitem{axioms} A. Jones, L. L. Littlejohn, A. Quintero Roba,
\emph{Krein-Sobolev Orthogonal Polynomials II}, Axioms 14 (2025), 115.
\url{https://doi.org/10.3390/axioms14020115}
\bibitem{lw1} L. L. Littlejohn, R. Wellman,
\emph{A general left-definite theory for certain self-adjoint operators with
applications to differential equations}, J. Differential Equations 181 (2002), 280--339.
\url{https://doi.org/10.1006/jdeq.2001.4079}
\end{thebibliography}
\end{document}
